\section{Cartographie mémoire \texttt{\%MW} des coupleurs NOC}
\label{secCartographieMemoire}

L'automate M340 de la cellule embarque deux coupleurs
\texttt{BMX~NOC~0401.2} (\ref{secEtapeConfigMaterielle}), dont la
réservation mémoire respective a été fixée comme suit :

\begin{itemize}
    \item le coupleur \texttt{NOC\_ETHIP\_CONVOYEUR}, qui ne scrute que le
    contrôleur PFC200 (\ref{secDonneesGuichet}), réserve \textbf{32 mots
    en entrée} à partir de \texttt{\%MW0}, puis \textbf{32 mots en
    sortie} juste après ;
    \item le coupleur \texttt{NOC\_ETHIP\_ROBOT}, dont l'image mémoire est
    \textbf{contiguë} à celle du premier coupleur, réserve \textbf{128 mots
    en entrée}, puis \textbf{128 mots en sortie}. Les équipements qu'il
    scrute ne sont pas détaillés ici.
\end{itemize}

\begin{UPSTIpreparation}[][Cartographie complète de l'espace \texttt{\%MW}]
    \UPSTIquestion{Établir la cartographie complète de l'espace
    \texttt{\%MW} occupé par les deux coupleurs : pour chacune des quatre
    zones (entrée et sortie de chaque coupleur), donner les bornes des
    zones \texttt{HEALTH\_BITS} et \texttt{CONTROL\_BITS} et celles de la zone
    image des données.}

    \UPSTIquestion{En déduire, par chevauchement (\emph{alias}) avec les
    mots de la zone image de données du coupleur
    \texttt{NOC\_ETHIP\_CONVOYEUR}, les adresses \texttt{\%MW.bit} exactes
    des quatre capteurs et des cinq actionneurs du poste guichet définis en
    \ref{secDonneesGuichet} (rappel : le PFC200 produit 12~octets et n'en
    consomme que 4, seul le premier octet de chaque zone étant
    significatif).}
\end{UPSTIpreparation}

\UPSTIremarque[Utilité de cette cartographie]{Cette cartographie sera
nécessaire pour vérifier, à l'\textbf{étape~5} de la procédure de
configuration (\ref{secEtapeVerification}), que l'éditeur de données de
Control Expert affiche bien les variables \texttt{guichet\_IN} et
\texttt{guichet\_OUT} aux adresses attendues.}
